\section{Weighted echoes, a resonance comb, and fixed-polynomial corrections}
\label{sec:weighted}
This section consolidates the weighted-observable, echo, resonance-comb,
and polynomial-correction notes. All statements use the fixed Pell map and
odd prime K, with $N+F\sqrt2=(1+\sqrt2)^K$, $D=(N-1)/2$,
$B_j=U_0^jBU_0^{-j}$, and $\tau_-=D^{-1}\operatorname{Tr}_-$.
The parameter $\tau$ below is a real kick amplitude, not a trace symbol.
The preceding normal form and lifting lemmas supply the stated inputs;
no independence of the finite Weyl operators is assumed.

\subsection{Bounded arrays and the ordered Gaussian law}
For real arrays f define
\[
 X_K(f)=(2K)^{-1/2}\sum_{j<K} f_j B_j,\qquad
 Q_K(f,g)=K^{-1}\sum_{j<K}f_jg_j.
\]
For fixed r and arrays uniformly bounded by M, the ordered mixed moments obey
\begin{equation}\label{eq:weighted-moments}
 \tau_-\bigl(X_K(f_1)\cdots X_K(f_{2r})\bigr)
 =\sum_{\pi\in\mathcal P_2(2r)}\prod_{\{a,b\}\in\pi}Q_K(f_a,f_b)
       +O_{r,M}(K^{-1})+O_{r,M}(K^r/N).
\end{equation}
For each fixed odd m, after the lifting threshold,
\[
 |\tau_-(X_K(f_1)\cdots X_K(f_m))|
       \leq M^m(2K)^{m/2}/(N-1).
\]
Thus convergence of a finite set of empirical Gram matrices gives the Wick
moments of a commuting real Gaussian vector with that covariance.

\begin{proof}
The exact parity-projected word trace is
$e_N(\phi)(N\mathbf1_{v=0}-1)/(N-1)$. It bounds the noncentral contribution
as displayed. The largest-gap norm lift and bounded-cluster argument confine
nonpaired words to $O_r(K^{r-1})$ choices. Odd central words are impossible:
a lifted polynomial is divisible by $z^2-6z+1$, so its value at one is
divisible by four and cannot be a sum of an odd number of signs.
For each matching of a balanced word, pair-label summation with the phase
replaced by one gives exactly the product of Q's. Repeated labels cost
$O_{r,M}(K^{-1})$. Its crossing phase is a product of $z_{a-b}^{\pm1}$.
Each individual difference is uniform, and $\sum_d|z_d-1|<9$; therefore
$|\prod u_e-1|\leq\sum_e|u_e-1|$ bounds its weighted mean error by
$9M^{2r}$ times the number of crossings divided by K. Distinct crossings
need not be independent. Summation proves the formula.
\end{proof}
The exact four-letter support theorem also gives
\begin{equation}\label{eq:weighted-comm}
 \|[X_K(f),X_K(g)]\|_{2,-}
 \leq\frac{6}{\sqrt K}\sqrt{\frac{2N}{N-1}}
                   \|f\|_\infty\|g\|_\infty,\qquad K\geq11.
\end{equation}
Indeed, the coefficient of $[B_t,B_s]$ is
$(f_tg_s-f_sg_t)/(2K)$, bounded by $\|f\|_\infty\|g\|_\infty/K$.
The four Weyl terms for different unordered pairs have disjoint support.
Their full-space squared norm is therefore at most
$4\|f\|_\infty^2\|g\|_\infty^2K^{-2}\sum_{j=1}^{K-1}(K-j)a_j^2$.
Use $a_{K-j}=a_j$, $a_j^2\leq2a_j$, and $\sum a_j<9$, then restrict the
positive quadratic form to the odd sector. This proves
\eqref{eq:weighted-comm}; it is not an operator-norm assertion.

\subsection{The interaction-picture echo}
Let $U_\varepsilon=e^{i\varepsilon B}U_0$,
$\varepsilon=\tau/\sqrt{2K}$ and $E_K(t)=U_\varepsilon^tU_0^{-t}$.
Write $t=mK+s$, $0\leq s<K$, including negative t by the usual floor
convention. Define $w_{t,K}(j)=m+\mathbf1_{j<s}$.
The centered coefficients $c_j=(K-1-2j)/(2K)$ satisfy the exact identity
\[
 c_j+t/K-c_{(j-t)\bmod K}=w_{t,K}(j).
\]
Inserting the finite normal form, telescoping t unitary factors, and combining
the three exponentials with \eqref{eq:weighted-comm} gives
\begin{equation}\label{eq:echo-finite}
 \|E_K(t)-e^{i\tau X_K(w_{t,K})}\|_{2,-}
       \leq (11C+2)\tau^2/\sqrt K\qquad(|t|\leq CK).
\end{equation}
To detail the constant, the original normal-form error contributes
$6|t/K|\tau^2/\sqrt K$. The comparison has three factors with arrays c,
the constant $t/K$, and the negative shifted c. The two Duhamel errors sum
to at most $c_N(3|t/K|+3/4)\tau^2/\sqrt K$, where
$c_N=\sqrt{2N/(N-1)}\leq\sqrt{7/3}$. This yields the displayed bound.
The same argument uses adjoints at negative times.

Let $\mathcal B(a)$, $0\leq a\leq1$, be one standard Brownian motion and,
for $\alpha=m+a$ with $0\leq a<1$, put
\[
 Y(\alpha)=m\mathcal B(1)+\mathcal B(a).
\]
For $\beta=n+b$ its covariance is
$mn+mb+na+\min(a,b)$. The arrays $w_{t,K}$ have precisely this limiting
Gram matrix when $t/K\to\alpha$. Therefore for every fixed finite collection
of integer exponents $r_l$ and times $t_l/K\to\alpha_l$,
\[
 \tau_-\left(\prod_l E_K(t_l)^{r_l}\right)
 \longrightarrow
 \exp\left[-\frac{\tau^2}{2}\sum_{l,j}r_lr_j
                     \operatorname{Cov}(Y(\alpha_l),Y(\alpha_j))\right].
\]
The order of the factors on the left is retained. Commutator estimates allow
them to be combined in the limit; there is no common eigenbasis at finite K.
In particular the phase variance of one echo is
$\tau^2[m^2+(2m+1)a]$. The same Brownian interval is reused on each cycle:
fresh independent noise would give a different variance. This is a tracial
finite-dimensional-distribution statement, not convergence of quantum
measurement trajectories in a path space.

\subsection{The full bounded recurrence window}
For $1\leq s<K$, the fixed-point equations for $W(v)U_0^s$ and
$RW(v)U_0^s$ have unique solutions, so
$|\tau_-(W(v)U_0^s)|\leq1/D$. Hence, with $X_K=X_K(1)$,
\[
 |\tau_-(X_K^lU_0^s)|\leq(2K)^{l/2}/D.
\]
For $|t|\leq CK$, $K\nmid t$, let $A_0=C|\tau|$. Taylor's integral
remainder of degree $2r-1$ and the normal form yield
\[
 |\tau_-(U_\varepsilon^t)|
 \leq\frac{6C\tau^2}{\sqrt K}
 +\frac1D\sum_{l<2r}\frac{(A_0\sqrt{2K})^l}{l!}
 +\frac{A_0^{2r}}{(2r)!}\tau_-(X_K^{2r}).
\]
First let K grow at fixed r, then r grow. The last limiting bound is
$A_0^{2r}/(2^rr!)\to0$. Thus nonresonant traces vanish uniformly in this
window, whereas $\tau_-(U_\varepsilon^{mK})\to e^{-m^2\tau^2/2}$.
If $e^{i\theta_l}$ are the actual eigenvalues, the empirical measure of
$(e^{i\theta_l},e^{iK\theta_l})$ converges to Haar measure times the wrapped
normal law: its Fourier coefficients are traces at $a+bK$.
This is a two-scale law of one eigenvalue, not independence of two levels.

\subsection{Minimal Pell relations and the first polynomial correction}
For an integer polynomial P, write $\|P\|_1$ for its coefficient norm.
If $P\ne0$ and $P(\lambda)=0$, then $\|P\|_1\geq8$, with equality only for
$\pm z^j(z^2-6z+1)$. Indeed $P(1)$ is divisible by four, so its norm is even;
if it were at most six the leading term would dominate since $\lambda>5$.
At norm eight the leading coefficient is one in absolute value. After a sign
change write $P=z^L+c z^{L-1}+R$, $\|R\|_1=7-|c|$. Evaluation gives
$|(6+c)\lambda-1|\leq7-|c|$, forcing $c=-6$ and then
$R=z^{L-2}$. Thus the first genuinely nonpaired central words have length eight.
The bounded-cluster argument improves their normalized contribution at any
fixed even order to $O_r(K^{-3})$.

For a fixed simple crossing graph, expand the phase product through first
order in $z_d-1$. Any two distinct edges form a forest, so their label
differences are independent uniform variables. The remainder is therefore
$O(K^{-2})$. A matching with c crossings contributes
$1+c(S_K-1)+O(K^{-2})$, $S_K=K^{-1}\sum_dz_d$. The total crossing count is
$A_r=\binom{2r}{4}(2r-5)!!$: choose four interlaced endpoints and match the
rest. Four repeated positions cause one excess matching; their six balanced
sign assignments contribute $-3A_r/(2K)+O(K^{-2})$ to the moment: the
matching-indexed sum counts each such word with one excess matching, so
the excess is subtracted. The crossings contribute
$A_r(S_K-1)=-(\gamma-3/2)A_r/K+O(K^{-2})$, and the two corrections add
to $-\gamma A_r/K$. Since
$S_K=1-(\gamma-3/2)/K+O((KN)^{-1})$, it follows that
\begin{equation}\label{eq:edgeworth-mom}
 \tau_-(X_K^{2r})=(2r-1)!!-
       \frac{\gamma}{K}\binom{2r}{4}(2r-5)!!
       +O_r(K^{-2})+O_r(K^r/N).
\end{equation}
Here
\[
 \gamma=\frac32+2\sum_{j\geq0}
 \left[1-\cos\left(\frac{4\pi}{\sqrt2(1+\sqrt2)}\lambda^{-j}\right)\right].
\]
For every fixed real polynomial f and standard Gaussian Z,
\[
 \tau_-(f(X_K))=\mathbb Ef(Z)-\frac{\gamma}{24K}\mathbb Ef^{(4)}(Z)
             +O_f(K^{-2})+O_f(K^{\deg f/2}/N).
\]
This is a fixed-polynomial expansion, not a distribution-function Edgeworth
theorem. Constants are not controlled as the polynomial degree grows.

For completeness, here is the stronger crossing-graph expansion, with the
notation it needs. For a fixed simple graph $G$ on the pair labels, let
the labels $t_v$ be independent and uniform in $\mathbf Z/K\mathbf Z$,
orient each edge, choose signs $\sigma_e\in\{-1,1\}$, and put
$Z_G=\prod_{e=(u,v)}z_{t_v-t_u}^{\sigma_e}$, where $z_d$ are the crossing
phases above; $\mathbb E$ is the average over the labels. Then
\[
 \mathbb E Z_G=S_K^{|E(G)|}+O_G(K^{-3}+N^{-1}).
\]
Indeed, expand $\prod_e(1+(z_e^{\sigma_e}-1))$. Every subset of at most
two edges is a forest, so its edge differences are independent and
uniform, and its expectation is exactly $(S_K-1)$ per edge, resp.\
$(S_K-1)^2$ per pair; this matches the binomial expansion of
$S_K^{|E(G)|}$ through second order, the remaining binomial terms being
$O_G(K^{-3})$ since $S_K-1=O(K^{-1})$. A subset of at least three edges
either contains a three-edge forest, whose expected absolute defect
product is at most $(9/K)^3$ by $\sum_d|z_d-1|<9$, or is a triangle. For
a triangle the bound is on the defect product, not on the unit-modulus
phases themselves: the three differences sum to zero modulo $K$, their
signed representatives therefore sum to $0$ or $\pm K$ (for $\pm K$ all
three have the same sign, since each is at most $(K-1)/2$ in absolute
value), in both cases
$\sum_i|\delta_i|\le K$, i.e.\ $\sum\ell\geq(K-3)/2$ for
$\ell(d)=(K-1)/2-|\delta|$; the phase-rate bound
$|z_d-1|\le A\lambda^{-\ell(d)}/(1-\lambda^{-K})$ with
$A=4\pi/(\sqrt2(1+\sqrt2))$ then gives a defect product
$O(\lambda^{-K/2})=O(N^{-1})$, uniformly in the labels, and if one
difference is zero the product is zero. All remaining defect factors are
bounded by two. This proves the stated refinement.

At order eight, the minimal relation is
$(3,-4)+(3,4)-6(1,0)=0$. Among its 56 orders the seven half-symplectic phases
$-12,-8,-4,0,4,8,12$ each occur eight times. Both signs and K time shifts
therefore give the isolated nonpaired central contribution
\[
 \frac{N}{N-1}K^{-3}
  [1+2\cos(8\pi/N)+2\cos(16\pi/N)+2\cos(24\pi/N)].
\]
It is $7/K^3+O((K^3N)^{-1})$, in addition to all balanced contributions;
it is not the whole eighth moment. The sufficient eight-letter lifting
threshold holds for odd primes $K\geq29$. Its remaining odd-parity trace
error is at most $16K^4/(N-1)$.

\subsection{Subcritical microscopic concentration without a Gaussian input}
For a circle spectrum define $a^2=D^{-1}\sum_l\operatorname{dist}
(\theta_l,K^{-1}\mathbb Z)^2$. Arcs of halfwidth $w<1/(2K)$ about the roots
contain at least $D(1-a^2/w^2)_+$ points. Counting gaps and applying
Cauchy--Schwarz to cells of length $x/D$ gives
\[
 \frac1D\#\{S_l>x\}\leq\frac{2a^2}{w^2}+\frac KD+\frac{2Kw}{x},
 \qquad
 R_D(x)\geq\frac{(1-a^2/w^2)_+^2}{K/D+2Kw/x}-1.
\]
An eigenvalue matching for normal matrices follows by decomposing their
squared-eigenbasis-overlap matrix into permutations. Chord distance bounds
circle distance, and the finite normal form gives
\[
 a\leq\frac{3\varepsilon^2+
            (|\varepsilon|/\pi)\sqrt{N/(N-1)}}{\sqrt K}.
\]
Consequently $|\varepsilon|\sqrt K\to0$ implies $Ka\to0$, collapse of the
gap law to $\delta_0$, and divergence of $R_D(x)$ for each fixed x. Set
$w=(Ka)^{2/3}/K$ for the gap estimate and $w=2a$ for the pair estimate;
the case $a=0$ is direct. These conclusions need no Gaussian limit and
refer to the original phases, not folded phases or isolated old eigenspaces.
